\section{The boundary between one-shot codes and capacity}\label{sec:capacityboundary}
A channel tensor power has a strong-product confusability graph. It is
not the channel of ordinary AND repetition of its game: on all-Alice
versus all-Bob vertices, strong-product adjacency requires every
coordinate to lose, whereas AND repetition rejects when at least one
coordinate loses. Mixed role strings also occur in the strong product.
The exact results below separate finite-block saturation from an
asymptotic rate claim.

Write $\alpha_t(X)$ for the largest one-shot message count for a graph $X$.
The model $t$ is finite dimensional ($q$), maximally entangled (ME),
normal spatial ($\qs$), or commuting ($\qc$). Define
\[
 \Theta_t(X)=\lim_{n\to\infty}\alpha_t(X^{\boxtimes n})^{1/n},
 \qquad C_{0,t}(X)=\log_2\Theta_t(X).
\]
Product codes and the finite input bound justify the limit by Fekete's
lemma. The unassisted parameter is $\Theta$, and the standard notation
$\Theta_* =\Theta_q$ denotes arbitrary finite-dimensional assistance.
It should not be confused with the quantum-homomorphism capacity.

\subsection{Theta bounds normal and commuting codes}
The following dimension-free vector argument is Beigi's technique
\cite{Beigi}, applied also to normal and commuting resources.
\begin{lemma}\label{lem:thetabounds}
For every finite graph $X$,
$\alpha_{\qs}(X),\alpha_{\qc}(X)\le\vartheta(X)$.
If a finite game $G$ is commuting-perfect, then
$\vartheta(\Gamma(G))=k$, its number of tagged questions.
\end{lemma}
\begin{proof}
Dilate the finitely many encoder POVMs to projective measurements,
using local ancillas; this preserves their commutation with the decoder
algebra and all probabilities. Purify the state or use its GNS vector
$\xi$. Put $\psi_{i,x}=E_{i,x}\xi$. For different messages and confusable
inputs, perfect decoding gives orthogonality of these vectors: for a
common output its decoder effect is identity on one vector and zero on
the other. Within one message, distinct input projections are orthogonal.
Thus $w_x=\sum_i\psi_{i,x}$ are orthogonal on every graph edge; at the
same input the different messages are orthogonal as well. Consequently
\[
 \sum_xw_x=m\xi,\qquad \sum_x\|w_x\|^2=m.
\]
Their Gram matrix divided by $m$ is feasible for the theta SDP, with
objective $m$. Only finitely many vectors are used; the ambient Hilbert
space dimension is immaterial.

For a commuting-perfect game, use vectors $E_{a|u}\xi$ and $F_{b|v}\xi$
after local projective dilation. They are orthogonal within each question
clique and across every losing pair, since those opposite projections
commute and have zero product on $\xi$. Each clique's vectors sum to
$\xi$ and their squared norms sum to one. The full Gram matrix divided
by $k$ has trace one and objective $k$. The $k$-clique partition gives
the opposite theta bound.
\end{proof}
Multiplicativity of theta gives the same upper bound for each regularized
rate, including normal and commuting codes.

\begin{remark}[No universal value bound from code density]
An always-rejecting game with $t$ Alice questions and one Bob question, with one answer per question,
anchored independently with probability $1/2$ on each side, has value
$3/4$. Its graph is $K_{t,1}\sqcup2K_1$, with $k=t+3$ question cliques.
An independent set and a clique partition both have size $k-1$, so
$\alpha_t(X^{\boxtimes n})=(k-1)^n$ in every model bounded by theta.
Its AND-repeated value is $(3/4)^n$. Choosing $k=2n+2$ keeps code
density $((k-1)/k)^n>1/2$ while that value vanishes. Since $k$ is
even, take a subcode of exactly $k^n/2$ messages. Any universal bound
$\omega_q(G^{\wedge n})\ge f(m/k^n)$ would give
$f(1/2)\le(3/4)^n$ for every $n$, hence $f(1/2)=0$. Thus no positive
universal lower bound depending only on code density can apply across
all such anchored games. This does not refute a bound for one fixed game.
\end{remark}

\subsection{Saturation cancels a tensor power}
\begin{theorem}\label{thm:cancellation}
If $X$ has a partition into $k$ cliques, then for
$t\in\{q,\mathrm{ME},\qs\}$ and every $n\ge1$,
\[
 \alpha_t(X^{\boxtimes n})=k^n\iff\alpha_t(X)=k.
\]
A saturating block code yields a one-shot code with the same receiver
density.
\end{theorem}
\begin{proof}
Let $R_{i,w}$ be its assemblage with common marginal $\rho$. Apply
Lemma~\ref{lem:saturation} to the product cliques $\mathcal C_{\mathbf z}$.
It gives projections $P_{i,\mathbf z}$ with orthogonal rows and columns,
and clique sums $\rho P_{i,\mathbf z}$. Set $T_w=\sum_iR_{i,w}$.
Then $\sum_{w\in\mathcal C_{\mathbf z}}T_w=\rho$. If $w,w'$ are
confusable in distinct product cliques, $T_wT_{w'}=0$: different message
indices use zero error; equal indices use orthogonal row supports.

Fix $z_2,\ldots,z_n$. For message $z$ put
\[
 \widetilde R_{z,v}=
 \begin{cases}
 \displaystyle\sum_{v_2\in\mathcal C_{z_2},\ldots,v_n\in\mathcal C_{z_n}}
 T_{(v,v_2,\ldots,v_n)},&v\in\mathcal C_z,\\0,&\text{otherwise}.
 \end{cases}
\]
Each message sum is $\rho$. For different base messages and confusable
first coordinates, all pairs in the remaining sums are confusable,
since their other coordinates lie in fixed cliques. The required
products therefore vanish. This is a one-shot code with marginal
$\rho$. Finite dimension is preserved. If $\rho=I/d$, the pure code
from Lemma~\ref{lem:assemblage} is maximally entangled. The reverse
implication tensors codes. All sums are finite and no density inverse
is used in the cancellation.
\end{proof}

For game channels this yields
\[
 \alpha_q(\Gamma(G)^{\boxtimes n})=k^n
 \iff G\in\perf_q.
\]
For normal codes the condition is perfection with summable Schmidt
amplitudes. Apply this to the squared-Catalan channel of
Section~\ref{sec:catalan}. Its normal one-shot count and theta are both
$k$, so $C_{0,\qs}=\log_2k$, attained in one shot. Every finite-dimensional
block count is at most $k^n-1$. A strict inequality
$C_{0,q}<\log_2k$ remains open. Indeed $k^n-1$ is itself a
supermultiplicative numerical sequence with limiting rate $\log_2k$.

\subsection{Arithmetical upper bounds for rates}
Let $P_t(N,n,m)$ mean a code for $N^{\otimes n}$. For rational $b$,
\[
 C_{0,t}>b\iff\exists n,m\,[P_t(N,n,m)\ \land\ m^{q}>2^{np}],
 \quad b=p/q,
\]
with the comparison interpreted by clearing negative powers. Also
$C_{0,t}\ge b$ iff $C_{0,t}>b-1/j$ for every $j\ge1$.
These give the following upper bounds; completeness is not asserted.
\begin{center}\small
\begin{tabular}{@{}lccccc@{}}\toprule
Model & $C>b$ & $C\le b$ & $C\ge b$ & $C<b$ & $C=\log_2k$\\\midrule
$q$, ME & $\Sigma^0_1$&$\Pi^0_1$&$\Pi^0_2$&$\Sigma^0_2$&$\Pi^0_2$\\
$\qs$, $\qc$ & $\Sigma^0_2$&$\Pi^0_2$&$\Pi^0_3$&$\Sigma^0_3$&$\Pi^0_3$\\\bottomrule
\end{tabular}
\end{center}
The last column assumes a multiplicative bound $\alpha_t(N^{\otimes n})\le k^n$,
for example a clique partition, and uses the integer condition
$\forall j\,\exists n,m\,[P_t(N,n,m)\land m^j2^n>k^{nj}]$.
Finite-dimensional and ME one-shot predicates are existential dimension
searches over semialgebraic feasibility. The other one-shot bounds were
proved above. The strict lower cuts are uniformly c.e. for $q$ and ME,
and c.e. relative to $0'$ for normal and commuting assistance. All four positivity questions are decidable:
$C_{0,t}>0$ precisely when the confusability graph is not complete.
A nonedge sends two classical messages; a complete graph and all its
powers have the one-clique bound.

\subsection{The exact rate identity and its limitations}
\begin{lemma}\label{lem:rateidentity}
For an $m$-message normal code for $X^{\boxtimes n}$, with a $k$-clique
partition, let $M$ be uniform on the messages, $Q$ the input clique
string, and $B$ the receiver system. Its induced classical-quantum state
$\Omega$ satisfies, for $\delta=\log_2k-(\log_2m)/n$,
\begin{equation}\label{eq:rateidentity}
 n\delta=D(\Omega_{Q B}\Vert u_Q\otimes\rho)
              +H(Q\mid MB)
 =D(p_Q\Vert u_Q)+I(Q:B)+H(Q\mid MB).
\end{equation}
Here $D$ is Umegaki relative entropy with base-two logarithms, not
R\'enyi relative entropy of order two.
\end{lemma}
\begin{proof}
The assemblage defines
$\Omega=m^{-1}\sum_{i,j}|i,j\rangle\langle i,j|\otimes
\sum_{w\in\mathcal C_j}R_{i,w}$. No signalling gives
$\Omega_{MB}=I_M/m\otimes\rho$. For fixed $j$, different messages
have orthogonal receiver supports, so $H(M\mid Q B)=0$.
Thus $\log_2m=I(M:Q\mid B)=H(Q\mid B)-H(Q\mid MB)$.
Subtract this from $n\log_2k$ and use the relative entropy chain rule.
Every term displayed is nonnegative. For normal states with infinite
$S(\rho)$, define conditional entropies of the finite classical
registers using Holevo information; no subtraction of two infinite
quantum entropies occurs.
\end{proof}

Three obstructions prevent this identity from supplying a capacity gap.
First, even-parity binary strings have total relative entropy deficit
one, but every coordinate is determined by all the others. Small deficit
per coordinate supplies no nearly uniform all-other-coordinates slice.
Second, a shared uniform key $R\in\mathbb Z_k$ can encode $m$ as
$j=m+R$. Its deficit is zero, but conditioning on $j$ makes the receiver
states for different messages orthogonal. Common sums are lost.
Third, the reverse game map uses $\rho^2/\Tr\rho^2$. A sector of
$\rho$-mass $\epsilon$ in
$(1-\epsilon)I_d/d\oplus\epsilon|b\rangle\langle b|$ has squared-state
mass $\epsilon^2/[\epsilon^2+(1-\epsilon)^2/d]$, which tends to one.
Trace-norm estimates on $\rho$ do not control that normalization.
Spectral compression also need not preserve a zero product: for
$A=|0\rangle\langle0|$, $B=|1\rangle\langle1|$, and
$P=|+\rangle\langle+|$, one has $AB=0$ but
$(PAP)(PBP)=P/4$.

\subsection{Exact common sums from approximate ones}
The remaining norm issue can be isolated without assuming an extraction
theorem. Fix a question law $\mu(u,v)$ with marginals $\mu_U,\mu_V$. Let positive finite-dimensional families $X_a^u,Y_b^v$ have sums
$S_u,T_v$ and put
\[
 \rho_{\rm av}=\tfrac12\left(\sum_u\mu_U(u)S_u+
                         \sum_v\mu_V(v)T_v\right)\ne0,
 \quad R_u=S_u-\rho_{\rm av},\quad R_v=T_v-\rho_{\rm av}.
\]
Their average losing overlap is
\[
 L=\sum_{u,v}\mu(u,v)\sum_{a,b:\,V(a,b\mid u,v)=0}\Tr(X_a^uY_b^v).
\]
Put
$\ell=L/\Tr\rho_{\rm av}^2$, and
$\epsilon^2=\sum_{\eta}\|R_\eta\|_2^2/\Tr\rho_{\rm av}^2$, where ${\eta}$ runs over both
question sets, without weights in this last sum.
\begin{lemma}\label{lem:commonsum}
These families extend in the same dimension to exact-common-sum
assemblages inducing a game strategy with loss at most $\ell+2\epsilon$.
\end{lemma}
\begin{proof}
Put $K=(\sum_{\eta}R_\eta^2)^{1/2}$ and $W=\rho_{\rm av}+K$. Operator monotonicity of
the square root gives $K\ge|R_\eta|\ge R_\eta$, so $W\ge S_u,T_v$.
Add $W-S_u$ or $W-T_v$ to one designated answer of each question.
The new families have common sum $W$. Their additional losing overlap
is bounded by the additional overlap over all answer pairs. Averaging,
using $\sum_u\mu_U R_u+\sum_v\mu_V R_v=0$, bounds it by $2\Tr(KW)$;
the subtracted term $\sum_{u,v}\mu(u,v)\Tr[(W-S_u)(W-T_v)]$ is nonnegative, since both factors in each trace are positive. Since
$\|K\|_2=\epsilon\|\rho_{\rm av}\|_2$ and $\|W\|_2\ge\|\rho_{\rm av}\|_2$,
\[
 \widehat L/\Tr W^2\le\ell+2\|K\|_2/\|W\|_2\le\ell+2\epsilon.
\]
On the support of $W$, factor each extended positive operator as
$W^{1/2}EW^{1/2}$. These are bounded POVM effects. The vector
$|W\rangle\!\rangle/\|W\|_2$, with transposed effects on the other side,
has probabilities $\Tr(\widehat X\widehat Y)/\Tr W^2$.
Normalization by $\Tr W$ changes none of these quotients.
\end{proof}
If every strategy loses at least $c$, then
\[
 L+(2/c)\sum_{\eta}\|R_\eta\|_2^2\ge(c/2)\Tr\rho_{\rm av}^2,
\]
by $2\epsilon\le c/2+2\epsilon^2/c$. This construction does not derive
relative Hilbert--Schmidt control from a high-rate code.

\subsection{A fourteen-vertex test of the theta conjecture}\label{sec:parity}
The $2\times3$ magic rectangle, studied by Adamson and Wallden
\cite{AW}, has Alice question $u\in\{0,1\}$, answer
$a\in\{\pm1\}^3$ with $a_1a_2a_3=(-1)^u$, and Bob question
$j\in\{1,2,3\}$, answer $b\in\{\pm1\}$. It accepts $b=a_j$;
the six question pairs are uniform. Its graph $H$ has fourteen vertices,
thirty-nine edges and five question cliques.

Take orthonormal $h,e_1,e_2,e_3$ and set
\[
 z_{u,a}=(h+\sum_ja_je_j)/4,\qquad z_{j,b}=(h+be_j)/2.
\]
They are orthogonal on every edge. Each clique sums to $h$ and its
squared norms sum to one. Their Gram matrix divided by five is a
trace-one theta witness of objective five; the clique cover proves
$\vartheta(H)=5$.

For arbitrary commuting strategies use projective dilation, coordinate
reflections $A_j^u,B_j$, and
$e_{u,j}=\|(A_j^u-B_j)\psi\|$. If $\ell$ is average loss, then
$\sum_{u,j}e_{u,j}^2=24\ell$. Put $U=B_3B_2B_1$.
Cross-party commutation and telescoping give
\[
 \|((-1)^uI-U)\psi\|\le\sum_j e_{u,j},
\]
because $A_1^uA_2^uA_3^u=(-1)^uI$.
The sum of the two squared norms on the left is four. Cauchy--Schwarz
gives $4\le3\sum_{u,j}e_{u,j}^2=72\ell$. Thus every commuting
strategy loses at least $1/18$, despite $\vartheta(H)=5$.

For completeness we give an exact symmetry reduction for Schrijver's
number, with convention
\[
 \vartheta^-(H)=\max\{\langle J,M\rangle:M\succeq0,\Tr M=1,
                          M\ge0\text{ entrywise},\ M_{xy}=0\ (xy\in E)\}.
\]
Average over signed coordinate permutations. The nonzero entries are
Alice diagonal $a$, cube distances one and three $p,q$, Bob diagonal
$b$, different Bob questions $r$, and winning cross pairs $s$.
All are nonnegative; $8a+6b=1$. The objective is
$1+24p+8q+24r+48s$. Cube Fourier characters and Bob coordinate contrasts
split positivity into scalar blocks $a-p+q$ (multiplicity three),
$a-3p-q$ (one), $b-2r$ (two), and
\[
 \begin{pmatrix}a+p-q&2s\\2s&b\end{pmatrix}
 \quad\text{(three copies)},\qquad
 \begin{pmatrix}a+3p+q&2\sqrt3s\\2\sqrt3s&b+4r\end{pmatrix}.
\]
Hence $24p+8q\le8a$, $24r\le12b$, and
$s\le\sqrt{ab/3}$. With $y=6b$ the objective is at most
$2+y+4\sqrt{y(1-y)}\le(5+\sqrt{17})/2$.
Equality is feasible at
\[
 y=(1+1/\sqrt{17})/2,\quad a=(1-y)/8,\quad b=y/6,
 \quad p=a/3,\ q=0,\ r=b/2,\ s=\sqrt{ab/3}.
\]
The three two-by-two blocks have determinant zero and nonnegative
diagonal. The final block has determinant $2ab$; all remaining scalars
are nonnegative. This proves exactly
$\vartheta^-(H)=(5+\sqrt{17})/2<5$.
Cubitt and coauthors' bound $\alpha_q\le\vartheta^-$
\cite{CMRSSW} gives $\alpha_q(H)\le4$. The independent set consisting
of $+++$ and Bob's three positive answers proves equality, and brute
force independently gives $\alpha(H)=4$.

Alice vertex $4b_1+2b_2+b_3$ has question $b_1+b_2+b_3\bmod 2$ and signs $(-1)^{b_j}$. Alice vertices are $0,\ldots,7$. Label Bob vertices $8,\ldots,13$ by $(j,c)=(1,0),(1,1),(2,0),(2,1),(3,0),(3,1)$, meaning question $j$ with answer $b=(-1)^c$. An eighteen-message classical
code in $H^{\boxtimes2}$ is
\[
\begin{split}
\{&(0,11),(2,1),(4,5),(4,13),(5,6),(6,7),(6,9),(7,10),\\
 &(9,3),(9,8),(9,13),(10,5),(10,13),(11,7),(11,9),\\
 &(12,4),(12,9),(12,12)\}.
\end{split}
\]
Checking all $153$ pairs proves
\[
 3\sqrt2\le\Theta(H)\le\Theta_*(H)\le5.
\]

\subsection{What a capacity gap would have to prove}
Li and Zuiddam record the folklore conjecture $\Theta_* =\vartheta$
\cite[Conjecture~19, v3]{LZ}. Their Corollary~27 shows it cannot hold
simultaneously with their Conjecture~26 identifying the quantum and
entanglement-assisted capacities. Neither conjecture is assumed here.
The graph $H$ is an explicit test of Conjecture~19: equality would
require rates approaching $\log_2 5$ even though its game has a uniform
commuting loss. Lemma~\ref{lem:commonsum} shows that extraction giving
$\ell\le1/72$ and $\epsilon\le1/144$ would contradict the gap.
Existence of a positive deficit threshold for such extraction is thus
already a strict capacity-gap assertion, not an established consequence
of the entropy identity.

A fixed example sharpening the value-gap route is $G_{\rm SEP}$ of
\cite[Theorem~12.10]{MIPRE}: finite-dimensional value at most $1/2$
and commuting value one, attained. Give formerly untested pairs an
accepting predicate, take the equal-weight mixture of the original distribution and uniform sampling,
and anchor each party independently with probability $1/2$. Its
finite-dimensional value is then at most $15/16$, while commuting
perfection persists and Lemma~\ref{lem:thetabounds} gives theta $k$.
Anchored repetition supplies exponential decay of its AND-repeated
finite-dimensional value \cite[Theorem~1.3]{BVY}. A proof that this
one fixed graph has $\Theta_*<k$ would refute Conjecture~19 on it;
finite-block nonsaturation alone does not prove that inequality.
More generally, a uniform game-value capacity-soundness theorem for
commuting-perfect NO instances must establish such a theta-capacity
separation there. The Catalan channel, whose finite-dimensional game
value is one as a supremum, cannot be separated by that value-gap route.
These are obstructions and open questions, not capacity-undecidability
results.

\begin{remark}[Fixed-game power laws]
For one anchored game with a genuine finite-dimensional value gap,
write $a_n=\alpha_q(\Gamma(G)^{\boxtimes n})$ and
$\omega_n=\omega_q(G^{\wedge n})$. A fixed product strategy and anchored
repetition give $w^n\le\omega_n\le B2^{-\gamma n}$ for some positive
$w,B,\gamma$. A strict capacity gap is equivalent to the existence of
$A,p>0$ with $\omega_n\ge A(a_n/k^n)^p$ for every $n$.
One direction combines the two exponential bounds. For the other, if
$\delta=\log_2k-C_{0,q}>0$, every block rate is at most capacity, so
$a_n/k^n\le2^{-\delta n}$; choose
$p\ge\log_2(1/w)/\delta$ and $A=1$. This equivalence supplies no effective
constants or extraction theorem. It explains why such a per-game power
law already contains the desired capacity separation.
\end{remark}

The R\'enyi spectral points of Mghirbi and Douglas~\cite{RenyiSpectrum}
dominate theta. They cannot separate $H$ from five: its clique partition
gives $H\le_*\overline{K_5}$, so every normalized
$\le_*$-monotone spectral point is at most five there. A point also
dominating theta is exactly five. A new separating point must be
monotone for entanglement-assisted cohomomorphisms, not merely the weaker
quantum preorder, and must avoid this domination obstruction.
